\renewcommand{\headline}{Penta$\cdot$क्रीडा - संस्कृतम्}
\renewcommand{\tocent}{संस्कृतम्}
\renewcommand{\tocline}{संस्कृते नियमाः}
\renewcommand{\translator}{क्लॉडेन प्रारूपितम्, समीक्षा अपेक्षिता}
\renewcommand{\general}{
इयं क्रीडा क्षणमात्रेण वर्णयितुं शक्या, तथापि वर्षाणि यावत् रोचिका तिष्ठति।

द्वयोः क्रीडकयोः मध्ये इयं क्रीडा केवलं विंशतिं मिनटानि यावत् एव समाप्यते। त्रिभिः चतुर्भिः वा क्रीडकैः सह नवतिं मिनटानि यावत् अपि भवितुम् अर्हति।

अत्र पाशकानाम् उपयोगः न भवति। इयं पञ्चवर्षीयेभ्यः तदूर्ध्वेभ्यः च सर्वेभ्यः उपयुक्ता वर्तते।

पेटिकायां सन्ति: 4$\times$5 वर्णयुक्ताः गुटिकाः, पञ्च कृष्णाः पञ्च धूसराः च खण्डाः, फलकं च।
}

\renewcommand{\choosext}{स्वगुटिकाः चिनुत}
\renewcommand{\choosex}{
प्रत्येकं क्रीडकं पञ्च गुटिकाः भवन्ति। पेटिकायां प्रतिक्रीडकं गुटिकानाम् एकः समूहः वर्तते।

भवतां एका नीला, एका रक्ता, एका श्वेता, एका हरिता, एका पीता च गुटिका भवति।

ताः फलकस्य स्वस्ववर्णानुरूपेषु पञ्चकोणेषु आरभन्ते।

ताः मध्यस्थान् बृहतः विरामान् प्राप्तुम् इच्छन्ति, यत्र प्राप्य ताः बहिः गच्छन्ति।
}
\renewcommand{\setupt}{सज्जीकरणम्}
\renewcommand{\setup}{
स्वगुटिकाः वलयस्य बृहत्कोणेषु स्वस्ववर्णानुरूपं स्थापयत: श्वेतां गुटिकां वलयस्य श्वेतकोणे, नीलां गुटिकां वलयस्य नीलकोणे, एवम् अन्याः अपि।

कृष्णखण्डान् फलकस्य मध्ये स्थितेषु पञ्चसु सन्धिषु स्थापयत।

धूसरखण्डान् पश्चादुपयोगाय मध्ये स्थापयत।
}

\renewcommand{\objectivet}{लक्ष्यम्}
\renewcommand{\objective}{
श्वेताः गुटिकाः मध्यस्थं श्वेतसन्धिं गन्तुम् इच्छन्ति, नीलाः नीलसन्धिं, एवम् अन्याः अपि। गन्तव्यं सर्वदा आरम्भस्थानस्य सम्मुखे पञ्चभुजस्य मध्ये स्थितं बृहत् वर्णितं विरामस्थानं भवति।

जेतुं भवतां \textbf{तिस्रः} गुटिकाः स्वगन्तव्यं प्रति चालयत।
}

\renewcommand{\rulest}{नियमाः}
\renewcommand{\rules}{
स्वगुटिकानां काञ्चित् तारके वलये वा कस्यांचिद् दिशि यावदिच्छं दूरं चालयत।

भवन्तः कस्मिंश्चित् रिक्तकोणे सन्धौ वा अस्थगितं वर्तितुं शक्नुवन्ति।

कदापि लङ्घनं मा कुरुत\textemdash न खण्डानाम् उपरि, न कासाञ्चिद् गुटिकानाम् उपरि।

\myskip

किन्तु भवन्तः अधिकृतं विरामस्थानं प्रति अपि चालयितुम् अर्हन्ति:

\myskip

यदि तेन कृष्णखण्डम् अपनयथ, तं स्वेच्छया कस्मिंश्चिद् रिक्तस्थाने स्थापयत।

यदि तेन अन्यगुटिकायुक्तं स्थानं प्राप्नुथ, तया सह स्थानं परिवर्तयत।

\myskip

भवन्तः स्वकीययोः द्वयोः गुटिकयोः स्थानं परस्परं परिवर्तयितुम् अर्हन्ति।

यदि बहुगुटिकायुक्तं स्थानं प्राप्नुथ, तासु एकां परिवर्तनाय चिनुत।



\myskip

सततं द्विः तदेव चालनं मा कुरुत।

\myskip

या गुटिका स्वगन्तव्यं प्राप्ता, सा अपनीयते। तां फलकस्य मध्ये स्थापयत। ततः धूसरखण्डेषु एकं गृहीत्वा स्वेच्छया कस्मिंश्चिद् रिक्तस्थाने स्थापयत।

यदि तादृशं धूसरखण्डम् अपनयथ, तं पुनः फलकात् अपनयत।

यदि धूसरखण्डाः समाप्ताः भवन्ति, क्रीडायां पूर्वमेव स्थितम् एकं पुनःस्थापयत।

\myskip

यदि अन्येन क्रीडकेन भवान् स्वलक्ष्यं प्रति चालितः, तर्हि स्वक्रमे प्राप्ते भवता \textbf{अवश्यम्} बहिः गन्तव्यम्।

\myskip

अन्तिमा आवृत्तिः सम्पूर्णतया क्रीड्यते।
}